\section{Introduction}\label{sec:introduction}

\noindent\textbf{Disclosure of AI assistance.}
The mathematical content and text of this manuscript were generated entirely by an autonomous exploration system driven by OpenAI's large language model Astra under human guidance.
No part of the work has been checked and understood by the author.
This manuscript records our exploration and its findings.
We share it in its current form to invite feedback and collaboration from interested readers.
\medskip

This work was motivated by a question posed by Roland Speicher in his blog \cite{Spe26}: whether the Fuglede--Kadison determinant formula for centered operator-valued semicircular elements extends to arbitrary mean.
It is also an experiment in doing mathematics with AI.
\medskip

The Fuglede--Kadison determinant extends the absolute determinant of a finite matrix to operators in a finite tracial von Neumann algebra.
For a matrix-valued semicircular element with zero mean, Mai and Speicher \cite[Theorem 1.2]{MS24} express this determinant through the capacity of a completely positive map determined by the coefficients.
A deterministic term changes the question substantially.
Such terms arise in particular when one studies the determinant of $T-zI$ as a function of the spectral parameter $z$.

In this paper we express the Fuglede--Kadison determinant of an affine semicircular pencil with complex matrix coefficients as an infimum over positive definite matrices.
We prove that the infimum is attained exactly when the pencil is invertible, and determine the unique minimizer in that case.

\subsection{Setting and the coefficient maps}
Let $(\M,\tau)$ be a finite von Neumann algebra with faithful normal tracial state.
Throughout, $s_1,\ldots,s_n\in\M$ are freely independent standard semicircular elements.
Each has spectral density $(2\pi)^{-1}\sqrt{4-t^2}$ on $[-2,2]$.
Freeness means that the trace of a product of centered elements from the generated unital algebras vanishes whenever adjacent factors come from different algebras;
see \cite[Section 1.11 and Chapter 2, Section 2]{MS17}.

For $m\ge1$, write $\Tr_m$ for the ordinary matrix trace, $\tr_m=\Tr_m/m$, $\varphi_m=\tr_m\otimes\tau$, and $E_m=\id_{M_m(\C)}\otimes\tau$.
Subscripts on $\Tr$ and $\varphi$ will be omitted when the size is fixed.
The notation $b>0$ means that the finite matrix $b$ is positive definite, while $b\ge0$ means positive semidefinite.
The identity matrix of size $m$ is written $1_m$, following \cite[Section 2.2]{MS24}.

\begin{definition}[Determinant and capacity]\label{def:determinant-capacity}
For $Y\in M_m(\M)$, let $\mu_{|Y|}$ be the spectral probability measure of $|Y|=(Y^*Y)^{1/2}$ under $\varphi_m$.
We use the Fuglede--Kadison determinant
\begin{equation}\label{eq:fk-definition}
 \Delta(Y)=\exp\left(\int_{[0,\infty)}\log t\dd\mu_{|Y|}(t)\right),
 \qquad \log0=-\infty,\quad \exp(-\infty)=0.
\end{equation}
For $b\in M_m(\C)$, put $\Delta_m(b)=|\det b|^{1/m}$.
For a completely positive map $\eta:M_m(\C)\to M_m(\C)$, its capacity is
\begin{equation}\label{eq:capacity-definition}
 \Cap(\eta)=\inf_{b>0}\frac{\det\eta(b)}{\det b}.
\end{equation}
\end{definition}
These are the conventions of \cite[Definitions 1.1 and 2.1]{MS24};
see also \cite{FK52,Gur04}.
Complete positivity means that every matrix amplification of the map preserves positivity.
In particular, maps of the form $b\mapsto\sum_j a_jba_j^*$ have this property \cite[Section 2.2]{MS24}.
The quotient in \eqref{eq:capacity-definition} is the homogeneous form of the definition with $\det b=1$.
Capacity also connects this determinant problem with operator scaling and noncommutative rank; see \cite[Sections 1.1--1.3]{GGOW20}.
The homogeneous function $\log\det\eta(b)-\log\det b$ is studied by geodesically convex optimization in \cite[Section 1.1]{AZGL+18}.

\begin{definition}[Affine semicircular pencil]\label{def:affine-pencil}
For $a_0,a_1,\ldots,a_n\in M_m(\C)$, an affine semicircular pencil is the operator
\begin{equation}\label{eq:affine-pencil}
 S=a_0+\sum_{j=1}^n a_j\otimes s_j\in M_m(\M).
\end{equation}
Here $a_0$ denotes $a_0\otimes1_{\M}$.
We use the same convention for other matrices in $M_m(\C)$.
We associate to the coefficients $a_1,\ldots,a_n$ the complex-linear maps
\begin{equation}\label{eq:coefficient-maps}
 \eta(b)=\sum_{j=1}^n a_jba_j^*,\qquad
 \rho(b)=\sum_{j=1}^n a_jba_j,
 \qquad b\in M_m(\C).
\end{equation}
\end{definition}
This terminology includes nonselfadjoint coefficients, as does the semicircular matrix convention in \cite[p.~3]{MS24}.
When we require selfadjoint coefficients, we state that assumption separately.
The maps $\eta$ and $\rho$ follow the notation of \cite[Theorem 1.2 and Lemma 3.5(ii)]{MS24}.

At $a_0=0$, the identity of Mai and Speicher is
\begin{equation}\label{eq:centered-capacity}
 \Delta\left(\sum_j a_j\otimes s_j\right)
       =e^{-1/2}\Cap(\eta)^{1/(2m)}.
\end{equation}
For a nonzero constant term, $a_0$ and $\eta$ alone do not determine the determinant.
Indeed, let $s$ be standard semicircular and $t>2$.
The scalar pencils $t+s$ and $t+is$ have the same constant term and the same map $\eta(b)=b$, but
\begin{equation}\label{eq:phase-obstruction}
 \Delta(t+s)<t<\Delta(t+is).
\end{equation}
Strict concavity of the logarithm gives $\tau\log(t+s)<\log t$, whereas $\tfrac12\tau\log(t^2+s^2)>\log t$.
Their maps $\rho$ are respectively $b\mapsto b$ and $b\mapsto-b$.
The general formula below retains this additional information.
Its appearance can also be seen directly in the covariance of the selfadjoint block operator associated with $S$.

\subsection{The general formula and attainment}
For each $b>0$, define
\begin{equation}\label{eq:inner-map}
 H_b=1_m+\eta(b),\qquad
 \mathcal T_b(Y)=H_bYb^{-1}+\rho(Y^*),\quad Y\in M_m(\C).
\end{equation}
This is a real-linear map on the real vector space of complex matrices.
Its invertibility for every $b>0$ is part of the next theorem.

% Provenance: Yin_C_F0126 and Yin_C_F0127; the manuscript proof uses the direct spectral-gap argument.
\begin{theorem}\label{thm:main}
Let $S$ be the pencil in Definition~\ref{def:affine-pencil}, with arbitrary complex coefficient matrices.
For every $b>0$, the map $\mathcal T_b$ in \eqref{eq:inner-map} is invertible as a real-linear map on $M_m(\C)$.
Write $Z_b=\mathcal T_b^{-1}(a_0)$.
Then
\begin{equation}\label{eq:main-determinant}
 \Delta(S)=e^{-1/2}\inf_{b>0}
 \left[\frac{\det(1_m+\eta(b))}{\det b}\right]^{1/(2m)}
 \exp\left[\frac1{2m}\operatorname{Re}\Tr(a_0^*Z_b)\right].
\end{equation}
Define
\begin{equation}\label{eq:main-action}
 \mathcal V(b)=\log\det H_b-\log\det b-m
                 +\operatorname{Re}\Tr(a_0^*Z_b).
\end{equation}
The following conditions are equivalent:
\begin{enumerate}
\renewcommand{\labelenumi}{(\roman{enumi})}
\item the infimum in \eqref{eq:main-determinant} is attained;
\item $\mathcal V$ has a critical point on the positive definite cone;
\item $S$ is invertible in the bounded operator algebra $M_m(\M)$.
\end{enumerate}
When these conditions hold, let $b_*$ denote the unique minimizer, which is also the unique critical point of $\mathcal V$, and put $Z_{b_*}=\mathcal T_{b_*}^{-1}(a_0)$.
Then
\begin{equation}\label{eq:main-optimizer}
 b_*=E_m((S^*S)^{-1}),\qquad Z_{b_*}=E_m((S^*)^{-1}).
\end{equation}
\end{theorem}

Equivalently,
\[
 \log\Delta(S)=\frac1{2m}\inf_{b>0}\mathcal V(b).
\]
The minimizers are the same as in \eqref{eq:main-determinant}.
Here a critical point means that the first variation of $\mathcal V$ vanishes in every selfadjoint matrix direction.

At zero constant term, $Z_b=0$ for every trial, and
\begin{equation}\label{eq:capacity-reduction}
 \inf_{b>0}\frac{\det(1_m+\eta(b))}{\det b}=\Cap(\eta).
\end{equation}
For the lower bound, use $1_m+\eta(b)\ge\eta(b)$ and determinant monotonicity.
Conversely, for every fixed $c>0$, the substitution $b=tc$ gives
\[
 \frac{\det(1_m+\eta(tc))}{\det(tc)}
   =\frac{\det(t^{-1}1_m+\eta(c))}{\det c}
   \longrightarrow\frac{\det\eta(c)}{\det c}.
\]
Taking the infimum over $c$ proves \eqref{eq:capacity-reduction}.
Hence our formula recovers \eqref{eq:centered-capacity} directly.
The proof of Theorem~\ref{thm:main} uses the semicircular Dyson equation, rather than the centered determinant identity, to establish the general formula.

\subsection{Other variational descriptions}
Throughout this subsection assume $a_j=a_j^*$ for $0\le j\le n$.
We first prove a strict-domain version of the formula proposed in Speicher's public discussion~\cite{Spe26}.
For $c>0$, put $\mathcal C_c(x)=cxc$ and
\begin{equation}\label{eq:hs-domain}
 \mathcal B_\eta=
 \{c>0:\mathcal C_c-\eta>0\text{ on }M_m(\C)\}.
\end{equation}
Here strict positivity is with respect to the Hilbert--Schmidt inner product $\langle x,y\rangle=\Tr(x^*y)$.

\begin{proposition}[Strict Hilbert--Schmidt formula]\label{prop:hs-formula}
For every $c\in\mathcal B_\eta$, the map $\mathcal C_c+\eta$ is invertible.
Put $h_c=(\mathcal C_c+\eta)^{-1}(a_0)$.
Then
\begin{equation}\label{eq:hs-determinant}
 \Delta(S)=e^{-1/2}\inf_{c\in\mathcal B_\eta}
       \Delta_m(c)\exp\left(\frac1{2m}\Tr(a_0h_c)\right).
\end{equation}
\end{proposition}
Section~\ref{sec:imaginary-hs} proves this formula and compares its trial domain and quadratic correction with those of Theorem~\ref{thm:main}.

Strictly positive pencils admit a minimum over a linear matrix inequality.
\begin{definition}[Feasible matrices]\label{def:positive-feasibility}
For $b>0$, put
\begin{equation}\label{eq:positive-lmi}
 L(b)=\begin{pmatrix}a_0-\eta(b)&1_m\\1_m&b\end{pmatrix}.
\end{equation}
The feasible and strictly feasible sets are respectively
\begin{equation}\label{eq:feasible-set}
 \mathcal B=\{b>0:L(b)\ge0\},\qquad
 \mathcal B_{\rm str}=\{b>0:L(b)>0\}.
\end{equation}
\end{definition}
Define
\begin{equation}\label{eq:positive-objective}
 J(b)=-\log\det b+\Tr(a_0b)-\tfrac12\Tr(b\eta(b))-m.
\end{equation}

% Provenance: the final matrix-supersolution formulation of Yin_C_F0121.
\begin{theorem}[Strictly positive pencils]\label{thm:positive}
The pencil $S$ is strictly positive, meaning $S\ge\gamma1_m$ for some $\gamma>0$, if and only if $\mathcal B_{\rm str}$ is nonempty.
Under this hypothesis,
\begin{equation}\label{eq:positive-minimum}
 \Delta(S)=\exp\left(\frac1m\min_{b\in\mathcal B}J(b)\right)
          =\exp\left(\frac1m\inf_{b\in\mathcal B_{\rm str}}J(b)\right).
\end{equation}
The minimum is attained uniquely at $q=E_m(S^{-1})$, which satisfies
\[
 q^{-1}=a_0-\eta(q),\qquad q\le b\quad(b\in\mathcal B).
\]
The infimum on $\mathcal B_{\rm str}$ is not attained.
\end{theorem}
Section~\ref{sec:positive} proves the theorem and relates $q$ to the minimizer $b_*=E_m(S^{-2})$ in Theorem~\ref{thm:main}.

Section~\ref{sec:ordered-minmax} gives ordered min--max formulas for scalar pencils, commuting coefficients and isotropic covariance, and a separate matrix formula when $\eta$ is strictly positive as a Hilbert--Schmidt operator.

\subsection{Consequences for Brown measure}
Let $T=a_0+\sum_j a_j\otimes s_j$.
The Brown measure $\mu_T$ is the probability measure whose logarithmic potential is $u(z)=\log\Delta(T-z1_m)$; the precise convention is recalled in Section~\ref{sec:spectral}, following \cite[Chapter 11, Section 9]{MS17}.
Spectral shifts in Theorem~\ref{thm:main} therefore give $u$ and the spectral criterion
\[
 z\notin\spec(T)
 \quad\Longleftrightarrow\quad
 \text{the infimum for $T-z1_m$ is attained}.
\]
Theorem~\ref{thm:main} and the nonnegative quadratic correction in \eqref{eq:main-action} give the following consequence.

% Proof input: Yin_C_F0172, Sections 1--2.
\begin{corollary}[Capacity bound and absence of atoms]\label{cor:brown-capacity}
For every $z\in\C$,
\begin{equation}\label{eq:intro-uniform-bound}
 \Delta(T-z1_m)\ge e^{-1/2}\Cap(\eta)^{1/(2m)}.
\end{equation}
If $\Cap(\eta)>0$, then $\mu_T$ has no atoms.
\end{corollary}
The proof is given in Section~\ref{subsec:brown-atoms}, where Proposition~\ref{prop:brown-atoms} also shows that every atom of $\mu_T$ lies in $\spec(a_0)$ without a capacity assumption.

When both $a_0$ and $\rho$ vanish, Theorem~\ref{thm:main} determines the radial distribution and the operator spectrum.
This includes linear combinations of free circular elements with matrix coefficients; see Example~\ref{ex:brown-circular-coefficients}.
Write $r(\eta)$ for the spectral radius of $\eta$, the maximum modulus of its eigenvalues.

% Proof input: Yin_C_F0170; positive capacity uses Yin_C_F0172.
\begin{proposition}[Radial Brown measure]\label{prop:brown-radial}
Assume $a_0=0$ and $\rho=0$.
The logarithmic potential has the form $u(z)=f(|z|^2)$, where $f$ is finite, continuously differentiable, nondecreasing and concave on $(0,\infty)$.
Its radial distribution is
\begin{equation}\label{eq:brown-radial-cdf}
 F(s):=\mu_T(\{|z|^2\le s\})=2sf'(s),\qquad s>0.
\end{equation}
The measure $\mu_T$ is absolutely continuous with respect to planar area on $\C\setminus\{0\}$.
If $\Cap(\eta)>0$, it is absolutely continuous on all of $\C$.
Moreover,
\begin{equation}\label{eq:brown-spectral-disk}
 \spec(T)=\{z\in\C:|z|^2\le r(\eta)\}.
\end{equation}
\end{proposition}
Section~\ref{subsec:brown-radial} proves this proposition and gives the distribution in terms of regularized minimizers, together with bounds for the density.
For general coefficients, Section~\ref{subsec:brown-density} computes the Laplacian of the regularized logarithmic potential through the Hessian at the minimizer.

\subsection{Proof mechanism and organization}
The proof of the general formula starts with a regularized variational formula.
The function being minimized has compact sublevel sets in the positive definite cone.
A positive-map argument shows that every critical point has the signs required by the selfadjoint block resolvent.
Dyson uniqueness then identifies the global minimizer.
Rescaling the trial variable removes the regularizer.
For the attainment assertion, a critical point yields a strict spectral-radius estimate and an analytic continuation of the matrix Cauchy transform through zero.
This proof is independent of the later Hilbert--Schmidt comparison and of the applications to Brown measure.

Section~\ref{sec:dyson} establishes the common Dyson and spectral tools.
Section~\ref{sec:main-proof} proves Theorem~\ref{thm:main}.
Section~\ref{sec:imaginary-hs} treats imaginary constant terms and the strict Hilbert--Schmidt comparison.
Section~\ref{sec:positive} gives the order-theoretic formula for strictly positive pencils, and Section~\ref{sec:ordered-minmax} establishes the ordered min--max principles.
Section~\ref{sec:spectral} develops the applications to spectra and Brown measures.
Appendix~\ref{app:comparison} contains the global comparison argument used in Section~\ref{sec:imaginary-hs}.
